% -----------------------------------------------------------------------------------
\section{Computable optimisation and its limits}\label{sec:computation}

This section asks what changes when the decision theory of Theorem~\ref{thm:universal} must be a
single algorithm reading descriptions. With a fixed substitution semantics and a termination promise,
nothing essential changes: one algorithm is optimal everywhere. Three kinds of obstruction appear
once the promise or the class is enlarged, and none of them concerns counterfactual inference.

\subsection{A uniform optimiser}

\begin{definition}[Promise class]\label{def:promise}
$\mathcal P_{\mathbb Q}$ is the class of program environments $E$ such that $E^\pi$ halts and returns an
exact rational value $V_E(\pi)$ for every $\pi\in\Pol_E$.
\end{definition}

The definition uses table oracles, so the validity of an environment is not defined in terms of
whichever optimiser is eventually constructed.

\begin{theorem}[Uniform optimisation on the promise class]\label{thm:promiseopt}
There is an agent that is promise-total on $\mathcal P_{\mathbb Q}$ and optimal under
$O_{\mathrm{oracle}}$ on every $E\in\mathcal P_{\mathbb Q}$.
\end{theorem}

\begin{proof}
On input $q(E,h)$, let $T^*$ recover $E$ and its finite policy space, run $E^\pi$ for every
$\pi\in\Pol_E$, select the first maximiser $\pi^*_E$ of $V_E$ in a fixed effective order, and return
$\pi^*_E(h)$. Every simulation halts by the promise, there are finitely many, and exact rational
comparison terminates. All owned queries of $E$ recover the same environment and use the same
tie-breaking, so the answers form the single policy $\pi^*_E$. Any agent $A$ induces the policy
$\pi^A_E(h)=A(q(E,h))$, and under oracle replacement
\[
  U(E,A)=V_E(\pi^A_E)\le\max_{\pi\in\Pol_E}V_E(\pi)=V_E(\pi^*_E)=U(E,T^*).\qedhere
\]
\end{proof}

$T^*$ never asks its live self what it would do during a simulated execution; it supplies table
answers instead. This is Theorem~\ref{thm:universal} made uniform: the lookup table is replaced by
an algorithm, and no knowledge of which parts of $E$ are predictors is needed, because oracle
substitution designates them. By Proposition~\ref{prop:progmatrix} the theorem maximises the
diagonal of $M_E$.

An explicit finite policy header is sufficient for the argument but not necessary. A finitely
branching oracle machine that halts for \emph{every} total oracle has a finite tree of consistent
oracle answers, since otherwise Kőnig's lemma \autocite{konig1927} gives an infinite path, whose
answers extend to a total oracle on which the machine runs forever; exhaustive exploration of the
tree then yields a finite optimisation problem. Halting against every \emph{computable} oracle does
not justify this compactness argument, since an infinite computable tree can lack a computable
infinite path.

Dependence of termination on the oracle is harmless as long as every policy terminates. The theorem
fails when some alternatives are undefined: if action $0$ returns $0$ and action $1$ simulates
machine $e$ and returns $1$ if it halts, the environment lies outside $\mathcal P_{\mathbb Q}$ unless
$e$ halts. Adding the convention that divergence pays $-1$ would make an exact optimiser choose $1$
exactly when $e$ halts, and so decide the halting problem. The obstruction then comes from the
added divergence semantics, not from the finite-policy argument.

\subsection{Two baseline impossibilities}

Two impossibilities are available with no counterfactual subtlety at all. Any result that claims
to be about counterfactual inference must avoid them.

\begin{proposition}[Foreign queries]\label{prop:foreign}
There are two terminating environments such that no agent, computable or not, is optimal on both
under $O_{\mathrm{oracle}}$.
\end{proposition}

\begin{proof}
Fix a binary query $q$ distinct from both environments' own descriptions, and let $E_0$ pay
$\one[A(q)=0]$ and $E_1$ pay $\one[A(q)=1]$. Each has optimal value $1$, attained by a constant agent,
and no agent attains $1$ in both.
\end{proof}

This is a cross-problem consistency contradiction requiring neither self-reference nor
undecidability; it is excluded by $\cls_{\mathrm{self}}$ with disjoint namespaces, and it
disappears under $O_{\mathrm{const}}$, which freezes the foreign response.

\begin{theorem}[The fixed-copy diagonal]\label{thm:fixedcopy}
If admissible agents are globally total and arbitrary fixed computable algorithms may occur in
environments, no agent is optimal on every environment, even among deterministic, binary,
rational-valued environments that consult the evaluated agent only on their own description.
\end{theorem}

\begin{proof}
Fix a globally total candidate $T$, and construct an environment $E_T$ containing the fixed algorithm
$g_0=T$. With $d=d(E_T)$, it computes $c=g_0(d)$, receives the evaluated agent's action $a$ at the
port, and pays $\one[a\neq c]$. The algorithm $g_0$ does not change when the evaluated agent changes.
Since $T$ is globally total, $g_0(d)$ halts, so $E_T$ terminates under either action. $T$ takes the
action $c$ and receives $0$; a globally total agent taking $1-c$ at $d$ receives $1$.
\end{proof}

In a programming formalism with no built-in parameter for the environment's own description,
Kleene's second recursion theorem \autocite{kleene1938} supplies the self-reference; see
\textcite{moschovakis2010} for a modern statement and proof. It produces a program with a specified
behaviour involving its own index; it does not prove that the behaviour terminates, and global
totality is what supplies termination. The proof holds under $O_{\mathrm{const}}$, and under
$O_{\mathrm{oracle}}$ whenever $g_0$ is an ordinary fixed program rather than a designated oracle
call. The transparent-and-opaque-box construction of \textcite{skalse2021} is the same mechanism with
different payoffs.

Remark~\ref{rem:skalse} found that construction ill defined in the model of
Section~\ref{sec:model}, and the contrast is instructive. There a problem was a finite object whose
data included the constant's output, so the construction required a problem $D_b$ with $b=T(D_b)$,
and an ex-ante optimiser has no such fixed point. Here the description of $E_T$ encodes a
computation involving $T$ without supplying its result, and no fixed point is needed: whatever $T$
outputs on $d(E_T)$, the copy outputs the same. Once that output $c$ is known, the environment is
the ordinary one-shot problem with payoffs $\one[a\ne c]$, and its maximiser is $1-c$. The agent
cannot uniformly select it while remaining globally total, because finding $c$ means running $T$ on
the very description it is reading.

There is no conflict with Theorem~\ref{thm:promiseopt}, because $T^*$ is only promise-total. Its
attempted diagonal environment contains a call to $T^*(d(E_{T^*}))$, whose evaluation invokes the
same optimisation again; the regress violates the table-termination promise, and the environment
is not in $\mathcal P_{\mathbb Q}$. The positive theorem supplies no globally total extension of $T^*$.

A third baseline is the program form of Death in Damascus.

\begin{proposition}[Self-prediction held fixed]\label{prop:doublecall}
Let an oracle environment call the agent twice on $d$, treating one answer as a prediction $p$ and the
other as the action $a$, and pay $\one[a\neq p]$. Under $O_{\mathrm{oracle}}$ every agent receives $0$
and every agent is optimal. Under $O_{\mathrm{const}}$ every agent has an improving deviation.
\end{proposition}

\begin{proof}
Coherent substitution answers both calls alike. An intervention changing only the action, with the
prediction frozen at the reference agent's answer $p$, earns $1$ with $a=1-p$.
\end{proof}

By Proposition~\ref{prop:progmatrix}, this is the matrix
$\bigl(\begin{smallmatrix}0&1\\1&0\end{smallmatrix}\bigr)$, in which no diagonal entry is maximal in
its row; it is Proposition~\ref{prop:did} again, a ratifiability obstruction and not a computability
result.

\subsection{Exact selection}

The following result avoids every baseline: there are no predictors, no copies of the agent, no
foreign queries, and the agent need halt only on valid inputs.

\begin{theorem}[No computable exact optimiser]\label{thm:exactsel}
No promise-total agent is an exact expected-payoff optimiser on all almost-surely terminating
binary-choice environments, even if realised payoffs lie in $\{0,1\}$, the agent is consulted only at
the action port on its own description, every expected value is uniformly computable, and every
environment has finite expected running time.
\end{theorem}

\begin{proof}
Let $(\varphi_e)$ be an effective enumeration of the partial computable functions, and let
$K_b=\{e:\varphi_e(e)\text{ halts with output }b\}$ for $b\in\{0,1\}$. The environment $E_e$ receives
the action $a\in\{0,1\}$, draws $N\in\mathbb N$ with $\Pr(N=n)=2^{-n-1}$, simulates $\varphi_e(e)$ for $N$
steps, and pays $1$ if the simulation halts within them with output $a$, and $0$ otherwise. Every
simulation is bounded, $E_e$ terminates almost surely, and its expected running time is finite. If
$e\in K_b$ with halting time $t$, then $V_e(b)=\Pr(N\ge t)=2^{-t}>0$ and $V_e(1-b)=0$; if neither
output occurs, both values are $0$.

Suppose a promise-total exact optimiser $T$ existed. Then $f(e)=T(d(E_e))$ is a total computable
function with $f(e)=b$ whenever $e\in K_b$. Let $e$ be an index of the total program $x\mapsto1-f(x)$.
Then $\varphi_e(e)=1-f(e)$, so $e\in K_{1-f(e)}$, and the condition requires $f(e)=1-f(e)$.
\end{proof}

The proof is the recursive inseparability of $K_0$ and $K_1$. What fails is \emph{exact selection},
not the approximation of expected values: after $n$ simulation steps the unresolved part of each
expectation is at most $2^{-n}$, so the values are uniformly computable reals with an explicit error
bound. What cannot be computed uniformly is which action belongs to the exact argmax when
arbitrarily small advantages and genuine ties are both admitted. The obstruction disappears as soon
as one of the two is excluded.

\begin{proposition}[Positive replacements]\label{prop:approxopt}
On a class of environments with finite policy spaces and uniformly computable values, there are
promise-total agents that are
\begin{enumerate}[label=(\roman*)]
  \item $\varepsilon$-optimal everywhere, for any fixed rational $\varepsilon>0$;
  \item exactly optimal on every environment with a unique maximising policy;
  \item exactly optimal on every environment whose values are exact rationals.
\end{enumerate}
\end{proposition}

\begin{proof}
(i) Approximate every value within $\varepsilon/3$ and maximise the approximations. (ii) Refine
enclosing intervals until one policy's lower bound exceeds every other policy's upper bound, which
happens in finite time when the maximiser is unique. (iii) is Theorem~\ref{thm:promiseopt}.
\end{proof}

Theorem~\ref{thm:exactsel} is therefore a true and non-trivial version of the thesis that no
computable rule is optimal on all fair problems, on an especially fair class. Its obstruction is
classical recursion theory, not a conflict among decision theories.

\subsection{Resource bounds}

Finite clocks or circuit descriptions can support globally total optimisers, because every required
evaluation then has a finite bound. They do not make optimisation as cheap as one execution.

\begin{proposition}\label{prop:sat}
If an agent running in time polynomial in the length of the description is exactly optimal on all
environments with $n$ binary information sets and a Boolean circuit as payoff, then
$\mathrm{P}=\mathrm{NP}$.
\end{proposition}

\begin{proof}
Given a Boolean formula with $n$ variables, form the environment whose information sets are the
variables and whose payoff is $1$ exactly on satisfying assignments. Query the agent on its $n$ owned
queries and evaluate the formula on the resulting policy. By optimality the policy is satisfying
whenever a satisfying assignment exists, so this decides satisfiability in polynomial time
\autocite{cook1971}.
\end{proof}

Per-input runtime bounds do not decide program equivalence either; the construction of
Theorem~\ref{thm:haltred} uses components bounded by $|x|$ steps on input $x$. The bounded variants
of Löb's theorem proved by \textcite{critch2019} yield positive cooperation results for specified
proof-responsive programs under proof-length assumptions; they do not establish optimality on a
class of environments of the present kind.

\subsection{What a fixed semantics settles}\label{sec:fixedsemantics}

Theorem~\ref{thm:promiseopt} extends at once to any environment class that supplies a
\emph{reference-independent} effective evaluator: a computable $V(E,\pi)$ on a finite policy space,
with exactly comparable values, whether the evaluator works by source analysis, by simulation, or by
proof. Finite circuits, explicitly clocked programs with stipulated time-out behaviour, and
descriptions carrying checkable certificates of the required evaluations all have such evaluators,
and admit optimisers that are globally total on the certified language, returning a default on
invalid strings. One might conclude that once a rule for propagating the dependence of the world on
the agent's output is fixed, optimal choice is a finite computation, and the only remaining question
is which rule.

The conclusion is false for anchored kernels. A fixed, computable kernel makes each comparison
computable, but it need not admit an agent whose actual output maximises it. In
Proposition~\ref{prop:doublecall} under $O_{\mathrm{const}}$ every baseline yields a finite,
unambiguous ranking, and every baseline is ranked below its alternative. Maximising a supplied table
$V_E$ is a finite computation; being a fixed point of a baseline-dependent kernel is a separate
problem, the ratifiability problem of Part~\ref{part:single} in program form.

Two further remarks fix the scope of \enquote{largest solvable class}. For reference-independent
problems with disjoint owned queries, every decidable class on which some globally total agent is
optimal can be enlarged by a single further problem, by altering the agent on that problem's owned
queries only; under $\cls_{\mathrm{self}}$ this changes nothing on the original class, so there is
no largest such class. Conversely, the class of all terminating, rational-valued, finite-policy
templates is a maximal promise domain for that representation, so an unqualified claim that there is
no largest solvable class would also be false.
